% GENERATED FILE. Edit canonical lessons, then run npm run generate:books.
% canonical-source-hash: fnv1a64:ace34cf6c460f460
% canonical-lessons: FR-C315-tombe, FR-C315-deuil, FR-C315-testament, FR-C315-heriter, FR-C315-heritage, FR-R315-revision-308-315

\chapter{Le deuil}
\label{ch:fr-le-deuil}
\hlchaptermodality{315}

\begin{chapteropening}
\textbf{By the end of this chapter:} \emph{I can talk about mourning and what people leave behind: the grave, mourning, the will, inheriting and an inheritance.}

Complete the last lesson of chapter 315: I can talk about mourning and what people leave behind: the grave, mourning, the will, inheriting and an inheritance.
\end{chapteropening}

\section[la tombe]{la tombe — the grave, the tomb}

\label{lesson:FR-C315-tombe}



\begin{quote}
Before the new one: say the French for the funeral, then the French for the cemetery.
\end{quote}

\subsection*{You'll want to know: la tombe}
\textbf{la tombe} — \textquotedblleft{}the grave, the tomb\textquotedblright{}.

\begin{quote}
Elle met des fleurs sur la tombe de sa mère.
\end{quote}

\begin{quote}
\emph{She puts flowers on her mother's grave.}
\end{quote}

Not to be confused with \emph{il tombe}, he falls.

\begin{cousinweb}
\emph{tombe} comes from Late Latin \emph{tumba} \textquotedblleft{}a tomb\textquotedblright{} (EtymDB), a borrowing from Greek \emph{týmbos} \textquotedblleft{}a burial mound\textquotedblright{}.
\end{cousinweb}

\begin{grammarlens}[title={What you've built}]
One more word for telling, explaining and describing at B1.
\end{grammarlens}

\subsection*{Your turn}
\begin{itemize}
\raggedright
  \item \emph{Say it:} \emph{la tombe}
  \item \emph{Say it:} \emph{la tombe}, once more
  \item \emph{Say it:} a sentence of your own with \emph{la tombe}
  \item \emph{Recall:} say the French for the funeral, then the French for the cemetery, then say \emph{la tombe} again
\end{itemize}

\begin{tcolorbox}[breakable,colback=teal!4,colframe=teal!35!black,title={Before you move on}]
What does la tombe mean? (\textquotedblleft{}The grave, the tomb\textquotedblright{}.) Say it once more.
\end{tcolorbox}
\section[le deuil]{le deuil — mourning, grief}

\label{lesson:FR-C315-deuil}



\begin{quote}
Before the new one: say the French for the cemetery, then the French for the grave.
\end{quote}

\subsection*{You'll want to know: le deuil}
\textbf{le deuil} — \textquotedblleft{}mourning, grief\textquotedblright{}.

\begin{quote}
Toute la famille est en deuil.
\end{quote}

\begin{quote}
\emph{The whole family is in mourning.}
\end{quote}

\emph{Porter le deuil} — to wear mourning; \emph{faire son deuil} — to come to terms with a loss.

\begin{cousinweb}
\emph{deuil} is usually traced to a Late Latin noun \emph{dolus} \textquotedblleft{}grief\textquotedblright{}, built on \emph{dolēre} \textquotedblleft{}to suffer, to grieve\textquotedblright{} (Lewis \& Short); it is not classical \emph{dolus} \textquotedblleft{}trickery\textquotedblright{}, and EtymDB lists no parent.
\end{cousinweb}

\begin{grammarlens}[title={What you've built}]
One more word for telling, explaining and describing at B1.
\end{grammarlens}

\subsection*{Your turn}
\begin{itemize}
\raggedright
  \item \emph{Say it:} \emph{le deuil}
  \item \emph{Say it:} \emph{le deuil}, once more
  \item \emph{Say it:} a sentence of your own with \emph{le deuil}
  \item \emph{Recall:} say the French for the cemetery, then the French for the grave, then say \emph{le deuil} again
\end{itemize}

\begin{tcolorbox}[breakable,colback=teal!4,colframe=teal!35!black,title={Before you move on}]
What does le deuil mean? (\textquotedblleft{}Mourning, grief\textquotedblright{}.) Say it once more.
\end{tcolorbox}
\section[le testament]{le testament — the will (legal)}

\label{lesson:FR-C315-testament}



\begin{quote}
Before the new one: say the French for the grave, then the French for mourning.
\end{quote}

\subsection*{You'll want to know: le testament}
\textbf{le testament} — \textquotedblleft{}the will (legal)\textquotedblright{}.

\begin{quote}
Il a laissé sa maison à ses enfants dans son testament.
\end{quote}

\begin{quote}
\emph{He left his house to his children in his will.}
\end{quote}

\emph{Faire son testament} — to make one's will.

\begin{cousinweb}
\emph{testament} is a learned borrowing from Latin \emph{testāmentum} \textquotedblleft{}a will\textquotedblright{}, from \emph{testārī} \textquotedblleft{}to bear witness\textquotedblright{}, built on \emph{testis} \textquotedblleft{}witness\textquotedblright{} (Lewis \& Short).
\end{cousinweb}

\begin{grammarlens}[title={What you've built}]
One more word for telling, explaining and describing at B1.
\end{grammarlens}

\subsection*{Your turn}
\begin{itemize}
\raggedright
  \item \emph{Say it:} \emph{le testament}
  \item \emph{Say it:} \emph{le testament}, once more
  \item \emph{Say it:} a sentence of your own with \emph{le testament}
  \item \emph{Recall:} say the French for the grave, then the French for mourning, then say \emph{le testament} again
\end{itemize}

\begin{tcolorbox}[breakable,colback=teal!4,colframe=teal!35!black,title={Before you move on}]
What does le testament mean? (\textquotedblleft{}The will (legal)\textquotedblright{}.) Say it once more.
\end{tcolorbox}
\section[hériter]{hériter — to inherit}

\label{lesson:FR-C315-heriter}



\begin{quote}
Before the new one: say the French for mourning, then the French for the will.
\end{quote}

\subsection*{You'll want to know: hériter}
\textbf{hériter} — \textquotedblleft{}to inherit\textquotedblright{}.

\begin{quote}
Elle a hérité de la maison de sa tante.
\end{quote}

\begin{quote}
\emph{She inherited her aunt's house.}
\end{quote}

\emph{Hériter de quelque chose} — with \emph{de}.

\begin{cousinweb}
\emph{hériter} comes from Old French \emph{eriter} (EtymDB), from Late Latin \emph{hērēditāre} \textquotedblleft{}to inherit\textquotedblright{}, built on \emph{hērēs} \textquotedblleft{}heir\textquotedblright{} (Lewis \& Short).
\end{cousinweb}

\begin{grammarlens}[title={What you've built}]
One more word for telling, explaining and describing at B1.
\end{grammarlens}

\subsection*{Your turn}
\begin{itemize}
\raggedright
  \item \emph{Say it:} \emph{hériter}
  \item \emph{Say it:} \emph{hériter}, once more
  \item \emph{Say it:} a sentence of your own with \emph{hériter}
  \item \emph{Recall:} say the French for mourning, then the French for the will, then say \emph{hériter} again
\end{itemize}

\begin{tcolorbox}[breakable,colback=teal!4,colframe=teal!35!black,title={Before you move on}]
What does hériter mean? (\textquotedblleft{}To inherit\textquotedblright{}.) Say it once more.
\end{tcolorbox}
\section[l'héritage]{l'héritage — the inheritance; the heritage}

\label{lesson:FR-C315-heritage}



\begin{quote}
Before the new one: say the French for the will, then the French for to inherit.
\end{quote}

\subsection*{You'll want to know: l'héritage}
\textbf{l'héritage} — \textquotedblleft{}the inheritance; the heritage\textquotedblright{}.

\begin{quote}
Il a reçu un petit héritage.
\end{quote}

\begin{quote}
\emph{He received a small inheritance.}
\end{quote}

Also a cultural legacy: \emph{l'héritage des anciens}.

\begin{cousinweb}
\emph{héritage} is built on \emph{hériter} \textquotedblleft{}to inherit\textquotedblright{} (EtymDB), from Late Latin \emph{hērēditāre}, built on \emph{hērēs} \textquotedblleft{}heir\textquotedblright{}.
\end{cousinweb}

\begin{grammarlens}[title={What you've built}]
One more word for telling, explaining and describing at B1.
\end{grammarlens}

\subsection*{Your turn}
\begin{itemize}
\raggedright
  \item \emph{Say it:} \emph{l'héritage}
  \item \emph{Say it:} \emph{l'héritage}, once more
  \item \emph{Say it:} a sentence of your own with \emph{l'héritage}
  \item \emph{Recall:} say the French for the will, then the French for to inherit, then say \emph{l'héritage} again
\end{itemize}

\begin{tcolorbox}[breakable,colback=teal!4,colframe=teal!35!black,title={Before you move on}]
What does l'héritage mean? (\textquotedblleft{}The inheritance; the heritage\textquotedblright{}.) Say it once more.
\end{tcolorbox}
\section[(dialogue)]{Retour sur les chapitres 308–315}

\label{lesson:FR-R315-revision-308-315}



\begin{quote}
Say the French for the birth and pregnant.
\end{quote}

\subsection*{Your turn}
Say these aloud:

\begin{itemize}
\raggedright
  \item the birth, pregnant, the delivery, to breastfeed, the cradle
  \item the maternity hospital, the christening, to christen, the pushchair, childhood
  \item adolescence, the teenager, youth, the start of the school year, of age
  \item to get engaged, the marriage proposal, to get married, the tiered wedding cake, the sugared almond
  \item the move, to move in, the house-warming, the tenant, the lease
  \item retirement, the retired person, old age, to take into hospital, the death
  \item to pass away, to bury, the burial, the funeral, the cemetery
  \item the grave, mourning, the will, to inherit, the inheritance
\end{itemize}

\begin{tcolorbox}[breakable,colback=teal!4,colframe=teal!35!black,title={Before you move on}]
The birth? (\textbf{la naissance}.) Pregnant? (\textbf{enceinte}.) The delivery? (\textbf{l'accouchement}.) To breastfeed? (\textbf{allaiter}.) The cradle? (\textbf{le berceau}.) The maternity hospital? (\textbf{la maternité}.) The christening? (\textbf{le baptême}.) To christen? (\textbf{baptiser}.) The pushchair? (\textbf{la poussette}.) Childhood? (\textbf{l'enfance}.) Adolescence? (\textbf{l'adolescence}.) The teenager? (\textbf{l'adolescent}.) Youth? (\textbf{la jeunesse}.) The start of the school year? (\textbf{la rentrée}.) Of age? (\textbf{majeur}.) To get engaged? (\textbf{fiancer}.) The marriage proposal? (\textbf{la demande en mariage}.) To get married? (\textbf{marier}.) The tiered wedding cake? (\textbf{la pièce montée}.) The sugared almond? (\textbf{la dragée}.) The move? (\textbf{le déménagement}.) To move in? (\textbf{emménager}.) The house-warming? (\textbf{la crémaillère}.) The tenant? (\textbf{le locataire}.) The lease? (\textbf{le bail}.) Retirement? (\textbf{la retraite}.) The retired person? (\textbf{le retraité}.) Old age? (\textbf{la vieillesse}.) To take into hospital? (\textbf{hospitaliser}.) The death? (\textbf{le décès}.) To pass away? (\textbf{décéder}.) To bury? (\textbf{enterrer}.) The burial? (\textbf{l'enterrement}.) The funeral? (\textbf{les funérailles}.) The cemetery? (\textbf{le cimetière}.) The grave? (\textbf{la tombe}.) Mourning? (\textbf{le deuil}.) The will? (\textbf{le testament}.) To inherit? (\textbf{hériter}.) The inheritance? (\textbf{l'héritage}.)
\end{tcolorbox}
